\subsubsection{Influence de (T, P) sur le potentiel chimique du dibrome}
\label{subsubsection:relation_potentiel_chimique_CP-TP}

\setcounter{equation}{0}
\setcounter{Q}{0}

On appelle $\mu^\star$ le potentiel chimique d'un corps monophasé, qui dépend de la température et de la pression.

\Q Déterminer la différentielle totale exacte de $\mu^\star (T, P)$.
\solution{
La différentielle totale exacte de G s'écrit : 
\begin{align*}
dG = -S dT + V dP + \mu^\star dn
\end{align*}
On peut alors en déduire le potentiel chimique du corps pur : 
\begin{align*}
d\mu^\star = -S_m dT + V_m dP
\end{align*}
}


On va chercher à représenter l'évolution du potentiel chimique à température $T$ constante.

\Q Tracer l'allure du potentiel chimique $\mu^\star (T, P^\circ) = \mu^{\star \circ} (T)$ du dibrome. Vous expliciterez les hypothèses faites.
\solution{
On cherche à comprendre le comportement en température du potentiel chimique, on a donc besoin de l'entropie molaire du dibrôme.
\begin{align*}
S^\star_m = - \frac{\partial \mu^\star}{\partial T}
\end{align*}
Si on suppose l'entropie molaire varie peu avec la température, on a : 
\begin{align*}
\int_{\mu^\star(T_0, P)}^{\mu^\star (T, P)} d\mu ' &= - \int_{T_0}^T S^\star_m dT' \\
\mu^\star (T, P) &= \mu^\star(T_0, P) - S^\star_m (T-T_0)
\end{align*}
Avant de tracer l'allure du potentiel chimique, on peut énumérer trois points : 
\begin{enumerate}
\item la variation de potentiel chimique est affine ; 
\item lors des transitions de phase, on a égalité de potentiels chimiques donc continuité de cette fonction ;
\item la pente change lors d'un changement de phase, avec $S^\star_{m,g} > S^\star_{m,l} >S^\star_{m,s}$
\end{enumerate}
On peut donc représenter l'évolution du potentiel chimique avec la température à une pression donnée $P^\circ$ : 
\begin{center}
\begin{minipage}{0.48\textwidth}
\includegraphics[width=\textwidth]{/home/jeff/Documents/IUTOrsay/1-Cours/4p13-Alliages/Cours-TD/Exos/CP/Trans-CP-Ex05/figures/Fig01}
\end{minipage}\begin{minipage}{0.48\textwidth}
\captionof{figure}{\'{E}volution du potentiel chimique du dibrome en fonction de la température}
\end{minipage}
\end{center}
Remarque : au vu de la formation des étudiants en thermodynamique, ceux-ci vont trouver l'utilisation d'une constante pour l'entropie molaire d'un gaz parfait contre intuitive et sortie tout droit du chapeau $\ldots$ Ils doivent l'admettre ... Ca peut se justifier à l'aide de la formule de Sackur-Tetrode (très loin de ce que nos étudiants connaissent) laquelle définie l'entropie d'un gaz parfait. Pour faire simple, elle stipule que l'entropie d'un gaz parfait évolue en température avec une relation de la forme : 
\begin{align*}
S_m = cste - \frac{3}{2} R \ln ( T )
\end{align*}
avec $R$ la constante des gaz parfaits.
Si la variation de température est modeste, l'entropie devrait variée de manière négligeable. Une variation de +\SI{10}{\celsius} conduit à une élevation de l'entropie molaire de 
\begin{align*}
\Delta S_m  &= S_m (342) - S_m(332) \\ 
&= -\frac{3}{2} \cdot 8.314 \cdot \ln \frac{342 + 273}{332 + 273} \\
&= -8.88 \cdot 10^{-2} ~ \si{\joule\per\kelvin\per\mole}
\end{align*}
La variation est relativement faible.
}

\Q En travaillant à une température fixée à \SI{298}{\kelvin}, donner la variation de $\mu^\star (\SI{298}{\kelvin}, P)$ lorsque la pression varie et que le corps pur évolue de la phase solide à la phase gazeuse. Vous expliciterez les hypothèses faites.
\solution{
On cherche à comprendre le comportement en pression du potentiel chimique, on a donc besoin du volume molaire du dibrome.
\begin{align*}
V^\star_m = \frac{\partial \mu^\star}{\partial P}
\end{align*}
Pour une phase condensée, on va supposer que la phase est incompressible, d'où $V^\star_m$ est constante avec une variation de la pression. On a : 
\begin{align*}
\int_{\mu^\star(T, P^\circ)}^{\mu^\star (T, P)} d\mu ' &= \int_{P^\circ}^P V^\star_m dP' \\
\mu^\star (T, P) &= \mu^\star(T, P^\circ) - V^\star_m (P-P^\circ)
\end{align*}
On a $\rho_s > \rho_l \implies V^\star_{m,s} < V^\star_{m,l}$, donc lors du passage d'un état solide à liquide, la pente augmente.
Pour un gaz, il faut exprimer le volume molaire d'un gaz parfait : 
\begin{align*}
V^\star_{m,g} = \frac{RT}{P}
\end{align*}
d'où on tire : 
\begin{align*}
\int_{\mu^\star(T, P^\circ)}^{\mu^\star (T, P)} d\mu ' &= RT \int_{P^\circ}^P \frac{1}{P'} dP' \\
\mu^\star (T, P) &= \mu^\star (T, P^\circ) + RT \ln \frac{P}{P^\circ}
\end{align*}
On peut donc représenter l'évolution du potentiel chimique avec la pression à une température donnée : 
\begin{center}
\begin{minipage}{0.48\textwidth}
\includegraphics[width=\textwidth]{\includegraphics[width=\textwidth]{/home/jeff/Documents/IUTOrsay/1-Cours/4p13-Alliages/Cours-TD/Exos/CP/Trans-CP-Ex05/figures/Fig02}
\end{minipage}\begin{minipage}{0.48\textwidth}
\captionof{figure}{\'{E}volution du potentiel chimique du dibrome en fonction de la pression à \SI{298}{\kelvin}}
\end{minipage}
\end{center}
}

\paragraph*{Données à \SI{298}{\kelvin}}
\begin{itemize}
\item Températures de transition de phase sous \SI{1}{\bar} : $T_{fus}^{\star \circ} = -7.2 ~ \si{\celsius}$ et $T_{vap}^{\star \circ} = 58.8 ~ \si{\celsius}$ ;
\item Entropies molaires du dibrome \ce{Br2} :
\begin{center}
\begin{tabular}{cccc}
\toprule
Etat physique & $\ce{Br2 (s)}$ & $\ce{Br2 (l)}$ & $\ce{Br2 (g)}$ \\
\midrule
$S_m^{\star \circ} ~ [\si{\joule\per\mole\per\kelvin}]$ & $112$ & $152$ & $245$\\
\bottomrule
\end{tabular}
\end{center}
\item Masses volumiques : 
\begin{center}
\begin{tabular}{ccc}
\toprule
Etat physique & $\ce{Br2 (s)}$ & $\ce{Br2 (l)}$ \\
\midrule
$\rho ~ [\si{\gram\per\centi\meter\cubed}]$ & $3.1$ & $3.1$\\
\bottomrule
\end{tabular}
\end{center}

Nota Bene : le dibrome solide est légèrement plus dense que le dibrome liquide, mais il faut plus de 2 chiffres significatifs.

\item \ce{Br2 (g)} est assimilé à un gaz parfait.
\end{itemize}